% =====================================================================
%  Rapport technique : maintenance prédictive d'un banc hydraulique
%  Compilation : pdflatex rapport.tex (deux fois, pour la table des matières)
%
%  - Les figures de rapport/figures/ et les tableaux de rapport/tables/ sont
%    produits automatiquement par le notebook (notebooks/projet_valve.ipynb).
%  - Les captures d'écran à faire soi-même vont dans rapport/captures/ avec le
%    nom indiqué dans chaque cadre « CAPTURE À INSÉRER ». Dès que le fichier
%    existe, il remplace automatiquement le cadre.
% =====================================================================
\documentclass[11pt,a4paper]{article}

\usepackage[utf8]{inputenc}
\usepackage[T1]{fontenc}
\usepackage[french]{babel}
\usepackage{lmodern}
\usepackage{microtype}
\usepackage[margin=2.3cm]{geometry}
\usepackage{graphicx}
\usepackage{booktabs}
\usepackage{tabularx}
\usepackage{multirow}
\usepackage{float}
\usepackage{caption}
\usepackage{subcaption}
\usepackage{xcolor}
\usepackage{enumitem}
\usepackage{fancyhdr}
\usepackage{amsmath,amssymb}
\usepackage{listings}
\usepackage[most]{tcolorbox}
\usepackage{csquotes}
\usepackage[colorlinks=true,linkcolor=blue!50!black,urlcolor=blue!50!black,citecolor=blue!50!black]{hyperref}

\graphicspath{{figures/}{captures/}}
\captionsetup{font=small,labelfont=bf}
\setlength{\parskip}{0.45em}
\setlength{\parindent}{0pt}

\definecolor{placeholder}{RGB}{214,83,64}
\definecolor{codebg}{RGB}{246,247,249}
\lstset{basicstyle=\ttfamily\small,backgroundcolor=\color{codebg},frame=single,rulecolor=\color{gray!40},
        breaklines=true,columns=fullflexible,keepspaces=true,showstringspaces=false,
        literate={é}{{\'e}}1 {è}{{\`e}}1 {à}{{\`a}}1 {ç}{{\c{c}}}1 {ê}{{\^e}}1 {É}{{\'E}}1 {ô}{{\^o}}1 {«}{{\guillemotleft}}1 {»}{{\guillemotright}}1}

% ---------------------------------------------------------------------
%  Résultats chiffrés exportés par le notebook (valeurs « ?? » sinon)
% ---------------------------------------------------------------------
\InputIfFileExists{tables/resultats.tex}{}{%
  \typeout{ATTENTION : tables/resultats.tex absent, exécutez le notebook.}}
\providecommand{\resnOutlierCycles}{??}\providecommand{\restrainPosPct}{??}
\providecommand{\restestPosPct}{??}\providecommand{\resnTest}{??}\providecommand{\resnTestPos}{??}
\providecommand{\resnFeatures}{??}\providecommand{\respctHighCorrPairs}{??}\providecommand{\resnAfterCorr}{??}
\providecommand{\resnLNonZero}{??}\providecommand{\reskDefault}{??}\providecommand{\resfinalModel}{??}
\providecommand{\resfinalCvAuc}{??}\providecommand{\resfinalCvFone}{??}\providecommand{\resfinalCvAcc}{??}
\providecommand{\restestAccuracy}{??}\providecommand{\restestPrecision}{??}\providecommand{\restestRecall}{??}
\providecommand{\restestFone}{??}\providecommand{\restestRocauc}{??}\providecommand{\restestTN}{??}
\providecommand{\restestFP}{??}\providecommand{\restestFN}{??}\providecommand{\restestTP}{??}
\providecommand{\resnTestErrors}{??}\providecommand{\resnFinalFeatures}{??}
\providecommand{\resriseHundred}{??}\providecommand{\resriseNinety}{??}\providecommand{\resriseEighty}{??}
\providecommand{\resriseSeventythree}{??}\providecommand{\resdelayNinetyMs}{??}
\providecommand{\resoneFeatTestAcc}{??}\providecommand{\resoneFeatTestFone}{??}

% ---------------------------------------------------------------------
%  Figures générées et captures d'écran (cadre de remplacement si absentes)
% ---------------------------------------------------------------------
% \placeholderbox{nom}{description de ce qu'il faut capturer}{hauteur}
\newcommand{\placeholderbox}[3]{%
  \begin{tcolorbox}[colback=placeholder!6,colframe=placeholder,arc=2mm,boxrule=0.9pt,
                    height=#3,valign=center,halign=center,width=0.92\linewidth]
    {\color{placeholder}\bfseries\large CAPTURE À INSÉRER}\\[0.4em]
    \small #2\\[0.5em]
    \footnotesize\ttfamily fichier attendu : rapport/captures/\detokenize{#1}.png
  \end{tcolorbox}}

% \capture{nom}{légende}{instruction}{hauteur}
\newcommand{\capture}[4]{%
  \begin{figure}[H]\centering
    \IfFileExists{captures/#1.png}{\includegraphics[width=0.92\linewidth,height=#4,keepaspectratio]{captures/#1.png}}%
      {\placeholderbox{#1}{#3}{#4}}
    \caption{#2}\label{fig:#1}
  \end{figure}}

% \figgen{nom}{légende}{largeur} : figure produite par le notebook
\newcommand{\figgen}[3]{%
  \begin{figure}[H]\centering
    \IfFileExists{figures/#1.png}{\includegraphics[width=#3\linewidth]{figures/#1.png}}%
      {\fbox{\parbox{0.8\linewidth}{\centering\vspace{1.5cm}Figure \texttt{figures/#1.png} :
       exécuter le notebook pour la générer.\vspace{1.5cm}}}}
    \caption{#2}\label{fig:#1}
  \end{figure}}

% \tabgen{fichier}{légende}{label}
\newcommand{\tabgen}[3]{%
  \begin{table}[H]\centering\small
    \caption{#2}\label{tab:#3}
    \IfFileExists{tables/#1.tex}{\resizebox{\ifdim\width>\linewidth\linewidth\else\width\fi}{!}{\input{tables/#1.tex}}}%
      {\fbox{Tableau \texttt{tables/#1.tex} : exécuter le notebook.}}
  \end{table}}

\pagestyle{fancy}\fancyhf{}
\lhead{\small Projet de Machine Learning}\rhead{\small Maintenance prédictive -- valve}
\cfoot{\thepage}
\renewcommand{\headrulewidth}{0.4pt}

% =====================================================================
\begin{document}

% ------------------------------ Page de garde ------------------------
\begin{titlepage}
  \centering
  {\large EFREI Paris -- ING2 BDML \par}
  {\large Machine Learning I -- Année universitaire 2026/2027\par}
  \vspace{3cm}
  {\rule{\linewidth}{0.6pt}\par\vspace{0.6cm}}
  {\Huge\bfseries Maintenance prédictive d'un banc hydraulique\par}
  \vspace{0.4cm}
  {\Large Prédiction de l'état optimal de la valve à partir des signaux de pression et de débit\par}
  \vspace{0.6cm}
  {\rule{\linewidth}{0.6pt}\par}
  \vspace{2.5cm}
  {\large Rapport technique\par}
  \vspace{1.5cm}
  {\large\textbf{[NOM Prénom -- à compléter]}\par
          \textbf{[NOM Prénom -- à compléter]}\par
          Groupe : \textbf{[à compléter]}\par}
  \vspace{1.5cm}
  {Enseignants : Ahmed Ghazi, Blaiech, Ahmad Tay, Hassan Shraim\par}
  \vfill
  {\large [Date de rendu -- à compléter]\par}
\end{titlepage}

% ------------------------------ Résumé -------------------------------
\begin{abstract}
Ce projet construit un système de maintenance prédictive pour un banc d'essai hydraulique :
pour chaque cycle de production de 60 secondes, il prédit si la valve fonctionne de façon
optimale (condition = 100\,\%) ou non. Les seules données utilisées sont la pression PS2
(100\,Hz), le débit FS1 (10\,Hz) et le fichier des conditions \texttt{profile.txt}. Après un
contrôle de qualité des données, chaque cycle est résumé par \resnFeatures{} caractéristiques
physiques (statistiques globales, dynamique, fenêtres temporelles, énergie spectrale), puis
une sélection de variables est apprise à l'intérieur de pipelines \texttt{scikit-learn}.
Neuf familles de modèles sont comparées en validation croisée stratifiée, optimisées par
Grid Search et Random Search, puis le modèle retenu (\resfinalModel) est évalué une seule fois
sur les \resnTest{} derniers cycles : accuracy \restestAccuracy, F1 \restestFone, AUC
\restestRocauc. L'AUC parfaite montre que le modèle classe correctement tous les cycles,
mais son seuil de décision est décalé par un changement de distribution entre
entraînement et test : toutes ses erreurs sont des valves à 90\,\% prédites optimales, alors
que les modèles non linéaires (SVM, forêt aléatoire, gradient boosting) atteignent 100\,\%
sur le test. L'analyse des variables identifie la cause physique d'une valve non optimale :
un \textbf{retard de la montée en pression} vers 9,4\,s dans le cycle (\resdelayNinetyMs\,ms dès
le niveau 90\,\%). Les bonus réalisés sont présentés (tests unitaires, conteneur Docker,
application web).
\end{abstract}

\tableofcontents
\newpage

% =====================================================================
\section{Introduction}

\subsection{Contexte et objectif}
L'usine souhaite détecter au plus tôt la dégradation de ses machines et comprendre pourquoi la
condition de la valve d'un cycle de production est parfois non optimale. On dispose pour cela
du jeu de données public \emph{Condition monitoring of hydraulic systems}
(UCI, Helwig et al., 2015), obtenu sur un banc d'essai composé d'un circuit principal de
travail et d'un circuit secondaire de refroidissement-filtration reliés par le réservoir
d'huile. Le banc répète des cycles de charge constants de 60\,s pendant lesquels l'état de
quatre composants (refroidisseur, valve, pompe, accumulateur) est volontairement dégradé.

\textbf{Objectif.} Construire un modèle de classification binaire qui prédit, pour chaque
cycle, si la condition de la valve est optimale (100\,\%) ou non. Les \textbf{2000 premiers
cycles} servent à construire le modèle ; les \textbf{\resnTest{} cycles restants} (2001 à 2205)
forment l'échantillon de test final.

\subsection{Données utilisées}
Conformément aux consignes, seuls trois fichiers sont utilisés (tableau~\ref{tab:capteurs}),
chaque ligne correspondant à un cycle.

\begin{table}[H]\centering\small
\caption{Fichiers utilisés.}\label{tab:capteurs}
\begin{tabular}{llllr}
\toprule
Fichier & Grandeur & Unité & Échantillonnage & Colonnes / cycle\\
\midrule
\texttt{PS2.txt} & Pression & bar & 100 Hz & 6000\\
\texttt{FS1.txt} & Débit volumique & l/min & 10 Hz & 600\\
\texttt{profile.txt} & Conditions des composants & -- & 1 par cycle & 5\\
\bottomrule
\end{tabular}
\end{table}

Les colonnes de \texttt{profile.txt} sont : condition du refroidisseur (3, 20, 100\,\%),
\textbf{condition de la valve} (100\,\% : commutation optimale ; 90\,\% : léger retard ;
80\,\% : retard important ; 73\,\% : proche de la panne), fuite interne de la pompe (0, 1, 2),
pression de l'accumulateur (130, 115, 100, 90 bar) et un indicateur de stabilité (0 : conditions
stables ; 1 : régime statique peut-être pas encore atteint).

\subsection{Organisation du code}
Le travail est entièrement reproductible :
\begin{itemize}[nosep]
  \item \texttt{notebooks/projet\_valve.ipynb} : notebook complet (toutes les étapes, figures et tableaux) ;
  \item \texttt{valve\_ml/} : module Python (chargement, extraction de caractéristiques, sélection, modèles),
        réutilisé par les tests et l'application ;
  \item \texttt{tests/}, \texttt{Dockerfile}, \texttt{app/} : bonus (section~\ref{sec:bonus}).
\end{itemize}
Les figures et tableaux de ce rapport sont exportés automatiquement par le notebook, ce qui
garantit que les chiffres cités correspondent exactement à l'exécution.

\capture{notebook_apercu}{Aperçu du notebook dans Jupyter.}
  {Capture du haut du notebook ouvert dans Jupyter/Colab (titre, plan et cellule d'imports).}{6cm}

% =====================================================================
\section{Prétraitement des données}

\subsection{Chargement et statistiques descriptives}
Les trois fichiers sont lus avec \texttt{pandas} (séparateur tabulation, sans en-tête). On
obtient deux matrices de signaux de dimensions $2205\times6000$ (PS2) et $2205\times600$ (FS1) et
un tableau $2205\times5$ pour les conditions. Le nombre de cycles est vérifié identique dans les
trois fichiers.

\capture{stats_descriptives}{Statistiques descriptives des signaux et des conditions (sortie du notebook).}
  {Capture des sorties des cellules « 1.2 Statistiques descriptives » : tableau PS2/FS1 et \texttt{profile.describe()}.}{6.5cm}

La pression PS2 varie fortement à l'intérieur d'un cycle (montées et descentes de charge),
alors que le débit FS1 présente des paliers entrecoupés de transitoires rapides
(figure~\ref{fig:exemples_cycles}). Les conditions des composants sont des niveaux discrets
fixés par le plan d'expérience.

\subsection{Valeurs manquantes et valeurs aberrantes}
Le contrôle qualité porte sur les valeurs manquantes, les valeurs infinies, les valeurs
physiquement impossibles (pression ou débit négatifs), les cycles constants (capteur figé) et
les lignes dupliquées.

\capture{qualite_donnees}{Contrôle de qualité des données (sortie du notebook).}
  {Capture du tableau « NaN / inf / valeurs < 0 / cycles constants / lignes dupliquées » et du tableau de détection des valeurs aberrantes (section 1.3).}{5.5cm}

Aucune valeur manquante ni infinie n'est présente : aucune imputation n'est nécessaire. Pour
les valeurs aberrantes, on calcule pour chaque cycle la moyenne, l'écart-type et le maximum de
chaque signal, puis on applique la règle de l'IQR ($x < Q_1 - 1{,}5\,\mathrm{IQR}$ ou
$x > Q_3 + 1{,}5\,\mathrm{IQR}$) et un z-score robuste (médiane et MAD, seuil 3,5).
\resnOutlierCycles{} cycles sont signalés par au moins un critère IQR
(figure~\ref{fig:boxplots_cycles_bruts}).

\figgen{boxplots_cycles_bruts}{Boxplots des statistiques par cycle utilisés pour détecter les valeurs aberrantes.}{1.0}

\textbf{Décision : les cycles atypiques sont conservés.} Ce ne sont pas des erreurs de mesure
mais des états physiques réels du banc (refroidisseur proche de la panne, accumulateur
dégradé, régime transitoire). Les supprimer retirerait précisément les situations qu'un
système de maintenance doit reconnaître, et de tels cycles existent dans le jeu de test.
Leur influence est maîtrisée par la standardisation et par l'utilisation de modèles robustes
(arbres et ensembles).

\subsection{Encodage des variables catégorielles}
Les capteurs ne fournissent que des variables continues. Les colonnes de \texttt{profile.txt}
sont des niveaux ordinaux déjà numériques ; elles ne sont pas utilisées comme variables
explicatives car ce sont des étiquettes d'état mesurées hors ligne, inconnues en production
(les utiliser constituerait une fuite d'information). Aucun encodage one-hot n'est donc
nécessaire ; la seule transformation est la binarisation de la cible.

\subsection{Transformation en classification binaire et équilibre des classes}
La cible est définie par
\[
  y = \begin{cases} 1 & \text{si la condition de la valve vaut } 100\,\% \text{ (optimale)},\\
                    0 & \text{sinon (90, 80 ou 73\,\%)}. \end{cases}
\]

\figgen{distribution_classes}{Distribution des niveaux de la valve et des classes binaires (entraînement / test).}{1.0}

Le jeu d'entraînement est quasiment équilibré (\restrainPosPct\,\% de cycles optimaux) : aucun
rééchantillonnage n'est nécessaire, mais la validation croisée est stratifiée et l'on suit le
F1-score et l'AUC en plus de l'accuracy. Le jeu de test est déséquilibré dans l'autre sens
(\restestPosPct\,\% de cycles optimaux, soit \resnTestPos{} sur \resnTest).

Le découpage imposé est \textbf{temporel}. Or les conditions des autres composants évoluent au
cours de l'expérience (figure~\ref{fig:conditions_temporelles}) : les 205 cycles de test ont
tous un refroidisseur à 100\,\% et un accumulateur majoritairement à 90 bar. Le jeu de test
n'est donc pas distribué comme le jeu d'entraînement (\emph{dataset shift}), ce qui rend
l'évaluation finale plus exigeante qu'une validation aléatoire.

\figgen{conditions_temporelles}{Évolution des conditions des composants au fil des cycles ; la ligne pointillée sépare entraînement et test.}{1.0}

\subsection{Visualisation des signaux}
La figure~\ref{fig:signaux_moyens_par_valve} superpose le profil moyen d'un cycle pour chaque
niveau de la valve. Le cycle comporte une phase sans pression (0 à 9\,s environ), une forte
montée en pression vers 9,4\,s accompagnée d'un pic de débit, puis des paliers de charge qui
changent toutes les 10\,s. Deux différences apparaissent nettement :
\begin{itemize}[nosep]
  \item la \textbf{montée en pression et le pic de débit arrivent d'autant plus tard que la
        valve est dégradée} : c'est la signature directe d'un retard de commutation ;
  \item sur les paliers suivants, la pression est légèrement plus haute et le débit plus bas
        quand la valve est dégradée (pertes de charge plus importantes).
\end{itemize}

\figgen{signaux_moyens_par_valve}{Profil moyen de PS2 et FS1 sur un cycle selon la condition de la valve.}{1.0}
\figgen{exemples_cycles}{Exemples de cycles bruts (vert : valve à 100\,\% ; rouge foncé : 73\,\%), cycle complet et zoom sur les 6 premières secondes.}{1.0}

\subsection{Extraction de caractéristiques}
Les 6600 points bruts d'un cycle sont très redondants et inadaptés à des modèles comme k-NN,
Naive Bayes ou un arbre unique. Chaque signal est donc résumé par des descripteurs
interprétables (tableau~\ref{tab:features}), soit \resnFeatures{} variables au total.

\begin{table}[H]\centering\small
\caption{Caractéristiques extraites pour chaque signal (PS2 et FS1).}\label{tab:features}
\begin{tabularx}{\linewidth}{lXr}
\toprule
Famille & Descripteurs & Nombre\\
\midrule
Statistiques globales & moyenne, écart-type, min, max, médiane, quantiles 5/25/75/95\,\%, IQR, amplitude, RMS, asymétrie, aplatissement, instants du max et du min & 16\\
Dynamique & moyenne de $|dx/dt|$, écart-type, max et min de la dérivée & 4\\
Fenêtres temporelles & moyenne et écart-type sur 10 fenêtres de 6\,s & 20\\
Spectre & part d'énergie dans 4 bandes logarithmiques, énergie totale (log) & 5\\
\midrule
Total par signal & & 45\\
\bottomrule
\end{tabularx}
\end{table}

\capture{code_features}{Code de l'extraction de caractéristiques (module \texttt{valve\_ml/features.py}).}
  {Capture de la cellule du notebook qui définit \texttt{extract\_signal\_features} et de la sortie « Matrice de caractéristiques : (2205, 90) ».}{6cm}

\subsection{Analyse des corrélations}
Toutes les analyses à partir de cette section sont faites sur les 2000 cycles d'entraînement
uniquement, pour ne pas laisser d'information du test influencer les choix.

\figgen{correlation_cible}{Les 15 variables les plus corrélées (Pearson) à la cible.}{0.75}
\figgen{matrice_correlation}{Matrice de corrélation des 20 variables les plus liées à la cible.}{0.8}

La matrice de corrélation montre de forts blocs de redondance : \respctHighCorrPairs\,\% des
paires de variables ont une corrélation absolue supérieure à 0,95 (par exemple moyenne,
médiane et RMS d'un même signal). Cette colinéarité gêne la régression logistique (coefficients
instables) et viole l'hypothèse d'indépendance de Naive Bayes, d'où le filtre de corrélation
de la section suivante. Les histogrammes et boxplots de la figure~\ref{fig:histogrammes_boxplots}
montrent que plusieurs variables séparent nettement les deux classes. Les plus corrélées à
la cible sont issues de la pression PS2 : moyenne sur la fenêtre 6--12\,s
(\texttt{PS2\_w1\_mean}, $r=0{,}77$), part d'énergie haute fréquence
(\texttt{PS2\_band3\_ratio}, $r=0{,}72$) et raideur maximale de la montée
(\texttt{PS2\_d\_max}, $r=0{,}70$). Toutes décrivent la montée en pression autour de 9,4\,s.
Les variables de débit FS1 sont moins corrélées ($|r|<0{,}5$) mais apportent une information
complémentaire.

\figgen{histogrammes_boxplots}{Histogrammes (haut) et boxplots (bas) par classe des six variables les plus discriminantes.}{1.0}

\paragraph{Lien avec les autres composants.} La figure~\ref{fig:valve_vs_autres_conditions}
et le V de Cramér (tableau~\ref{tab:cramer}) mesurent l'association entre la condition de la
valve et celles des autres composants.

\begin{table}[H]\centering\small
\caption{Association entre la condition de la valve et les autres variables de \texttt{profile.txt} (V de Cramér, 2205 cycles).}\label{tab:cramer}
\begin{tabular}{lcccc}
\toprule
& Refroidisseur & Pompe & Accumulateur & Indicateur de stabilité\\
\midrule
V de Cramér & 0,004 & 0,307 & 0,171 & 0,708\\
\bottomrule
\end{tabular}
\end{table}

Le refroidisseur est indépendant de la valve. Les associations avec la pompe,
l'accumulateur et surtout l'indicateur de stabilité viennent du \textbf{plan d'expérience} et
non d'un lien physique. En effet, chaque niveau dégradé de la valve (73, 80, 90\,\%) est testé
exactement 120 fois pour chaque état de la pompe et 90 fois pour chaque état de
l'accumulateur ; de plus, les 756 cycles « non stables » (phases de transition après un
changement de configuration) ont tous été enregistrés avec une valve à 100\,\%. Ces variables
ne sont donc pas des causes de la dégradation de la valve, et elles ne seraient de toute
façon pas disponibles en production. La dégradation doit être détectée dans la
\textbf{signature des signaux} de pression et de débit.

\figgen{valve_vs_autres_conditions}{Proportion de cycles avec valve optimale selon l'état des autres composants.}{1.0}

\subsection{Sélection de variables}
Quatre familles de méthodes sont combinées :
\begin{enumerate}[nosep]
  \item \textbf{filtre de variance} : suppression des variables constantes ;
  \item \textbf{filtre de corrélation} : pour chaque paire avec $|r|>0{,}95$ une seule variable est
        gardée (\resnFeatures{} $\rightarrow$ \resnAfterCorr{} variables) ;
  \item \textbf{méthodes univariées} : test ANOVA F et information mutuelle ;
  \item \textbf{méthodes embarquées} : régression logistique pénalisée L1 (\resnLNonZero{}
        coefficients non nuls pour $C=0{,}1$) et importance des variables d'une forêt aléatoire.
\end{enumerate}

\figgen{selection_variables}{Classement des variables selon l'ANOVA F, l'information mutuelle, la pénalisation L1 et la forêt aléatoire.}{1.0}

Les méthodes s'accordent sur les deux variables les plus utiles, \texttt{PS2\_w1\_mean} et
\texttt{PS2\_band3\_ratio} (meilleurs rangs moyens), ce qui renforce la confiance dans la
sélection. Elles diffèrent ensuite : l'information mutuelle et la forêt aléatoire valorisent
des variables à relation non monotone avec la cible (asymétrie, instant du maximum), que le
test ANOVA F, qui ne mesure qu'un écart de moyennes, classe plus bas. La pénalisation L1 ne
garde que \resnLNonZero{} variables, ce qui confirme que l'information est concentrée dans
peu de descripteurs. Pour la modélisation, le filtre de corrélation et la
sélection \texttt{SelectKBest} (ANOVA F) sont placés \textbf{à l'intérieur du pipeline} : ils sont
réappris dans chaque pli de validation croisée, ce qui évite toute fuite d'information. Le
nombre de variables $k$ est choisi en traçant l'AUC en validation croisée en fonction de $k$
(figure~\ref{fig:choix_k}) : on retient par défaut le plus petit $k$ dont l'AUC est à moins de
0,005 du maximum, soit $k=\reskDefault$ ; au-delà, l'AUC est saturée (proche de 1). Ensuite $k$ est de nouveau optimisé avec les
hyperparamètres (section~\ref{sec:optim}).

\figgen{choix_k}{AUC et F1 en validation croisée selon le nombre de variables retenues.}{1.0}

% =====================================================================
\section{Modélisation}

\subsection{Pipeline commun}
Tous les modèles partagent le même enchaînement, ce qui rend la comparaison équitable :
\begin{center}
\texttt{CorrelationFilter(0,95)} $\rightarrow$ \texttt{StandardScaler} $\rightarrow$
\texttt{SelectKBest(ANOVA F, k)} $\rightarrow$ classifieur.
\end{center}
La standardisation est indispensable pour k-NN, le SVM et la régression logistique ; elle est
neutre pour les arbres.

\subsection{Choix et justification des modèles}
Le tableau~\ref{tab:justif} résume les propriétés théoriques des modèles comparés et leur
adéquation aux données : variables continues, en partie corrélées, avec un bruit de mesure
et des cycles atypiques.

\begin{table}[H]\centering\small
\caption{Modèles comparés et justification.}\label{tab:justif}
\begin{tabularx}{\linewidth}{>{\raggedright}p{2.6cm}>{\raggedright}p{2.6cm}X}
\toprule
Modèle & Nature & Adaptation aux données\\
\midrule
Régression logistique & linéaire, paramétrique & Référence simple et interprétable. Suffisante si les classes sont séparables linéairement dans l'espace des caractéristiques ; sensible à la colinéarité (traitée par le filtre de corrélation et la régularisation L1/L2).\\
k-NN & non linéaire, à base d'instances & Frontière locale très flexible, sans hypothèse de distribution ; sensible aux échelles (standardisation), au bruit et à la dimension (d'où la sélection de variables).\\
Arbre de décision & non linéaire, non paramétrique & Insensible aux échelles et aux valeurs extrêmes, règles lisibles ; variance élevée (sur-apprentissage) si la profondeur n'est pas limitée.\\
Naive Bayes gaussien & probabiliste, génératif & Très rapide ; suppose des variables indépendantes et gaussiennes, hypothèses violées ici (variables corrélées, distributions multimodales) : sert de référence à biais élevé.\\
SVM (noyau RBF) & non linéaire, marge maximale & Bien adapté aux variables continues standardisées de dimension moyenne ; la marge le rend robuste au bruit ; demande le réglage de $C$ et $\gamma$.\\
Random Forest & ensemble (bagging) & Moyenne d'arbres décorrélés : réduit la variance, robuste au bruit, aux valeurs aberrantes et aux variables corrélées.\\
Gradient Boosting, AdaBoost & ensemble (boosting) & Corrige séquentiellement les erreurs : réduit le biais ; doit être régularisé (taux d'apprentissage, profondeur, sous-échantillonnage).\\
Stacking & ensemble hétérogène & Combine des modèles de biais différents (linéaire, noyau, arbres) par un méta-modèle logistique.\\
\bottomrule
\end{tabularx}
\end{table}

% =====================================================================
\section{Évaluation des modèles}

\subsection{Protocole}
Les modèles sont évalués par \textbf{validation croisée stratifiée à 5 plis} (mélangés,
graine 42) sur les 2000 cycles d'entraînement. Pour chaque modèle on mesure l'accuracy, la
précision, le rappel, le F1-score (classe positive : valve optimale) et l'aire sous la courbe
ROC (AUC). Les scores sur les plis d'entraînement sont aussi relevés : l'écart
entraînement--validation mesure la variance, le niveau du score de validation mesure le biais.

\subsection{Comparaison des modèles avant optimisation}
\tabgen{cv_base}{Validation croisée (5 plis) des modèles avec leurs hyperparamètres par défaut. « F1 train » est le F1 sur les plis d'entraînement, « Écart F1 » l'écart entraînement--validation.}{cv_base}
\figgen{comparaison_modeles_base}{Comparaison des métriques en validation croisée avant optimisation.}{1.0}

Le problème est très bien séparé dans l'espace des caractéristiques : tous les modèles sauf
Naive Bayes dépassent 0,98 de F1 et 0,98 d'AUC. Les ensembles (forêt aléatoire, gradient
boosting) et le SVM sont les meilleurs (F1 $\approx$ 0,995--0,997). La régression logistique,
modèle linéaire, atteint déjà 0,99 de F1 : la frontière entre les classes est presque
linéaire. Naive Bayes reste nettement en dessous (F1 $\approx$ 0,91) car ses hypothèses
d'indépendance et de normalité sont violées. L'arbre de décision non élagué a l'écart
entraînement--validation le plus grand (0,013), signe de variance.

\subsection{Matrices de confusion et courbes ROC}
Les prédictions hors-pli (\texttt{cross\_val\_predict}) permettent de construire, pour chaque
modèle, une matrice de confusion et une courbe ROC sur l'ensemble des 2000 cycles
d'entraînement, chaque cycle étant prédit par un modèle qui ne l'a pas vu.

\figgen{matrices_confusion_cv}{Matrices de confusion en validation croisée (seuil 0,5).}{1.0}
\figgen{roc_cv}{Courbes ROC en validation croisée.}{0.62}

\subsection{Compromis biais / variance}
Les courbes d'apprentissage (figure~\ref{fig:courbes_apprentissage}) se lisent ainsi : un modèle
à forte \textbf{variance} (arbre non élagué, k-NN avec peu de voisins) a un score
d'entraînement proche de 1 et un écart important avec la validation ; un modèle à fort
\textbf{biais} (Naive Bayes) a des courbes d'entraînement et de validation basses et proches.
Le bon compromis est obtenu par les modèles dont la courbe de validation est haute et proche
de la courbe d'entraînement.

Ici, l'arbre de décision et la forêt aléatoire ont un score d'entraînement de 1 quelle que
soit la taille du jeu : ils mémorisent les données. Pour l'arbre, l'écart reste visible même
avec 1600 cycles (variance élevée) ; la forêt réduit cet écart grâce au bagging. Naive Bayes
présente deux courbes basses et confondues autour de 0,90 : c'est un biais élevé, que plus
de données ne corrigera pas. La régression logistique, le SVM et k-NN ont des courbes qui
convergent vers 0,99--1 avec un faible écart : c'est le meilleur compromis biais/variance.

\figgen{courbes_apprentissage}{Courbes d'apprentissage (F1) : entraînement en rouge, validation en vert.}{1.0}

\subsection{Robustesse du protocole de validation}
Les cycles successifs partagent les mêmes conditions de fonctionnement, donc une validation à
plis mélangés place dans le pli de validation des cycles très semblables à ceux de
l'entraînement. Le tableau~\ref{tab:robustesse} compare cette validation à une validation à
\textbf{plis contigus} (sans mélange), plus proche de la situation réelle où l'on prédit des
cycles futurs.

\tabgen{robustesse}{Validation à plis mélangés vs plis contigus.}{robustesse}

Le résultat est net : avec des plis contigus, le F1 des ensembles d'arbres chute (0,997
$\rightarrow$ 0,935 pour la forêt aléatoire, 0,997 $\rightarrow$ 0,925 pour le gradient
boosting) alors que la régression logistique et le SVM restent vers 0,98. La validation
mélangée est donc optimiste, et les performances sur une période non vue dépendent de la
capacité du modèle à extrapoler à des conditions de fonctionnement nouvelles.

% =====================================================================
\section{Optimisation des hyperparamètres}\label{sec:optim}

\subsection{Méthodes}
\begin{itemize}[nosep]
  \item \textbf{Grid Search} (recherche exhaustive) pour les modèles à espace de recherche
        réduit : régression logistique ($C$, pénalité L1/L2), k-NN (nombre de voisins,
        pondération, distance), arbre de décision (profondeur, taille minimale des feuilles,
        critère), Naive Bayes (lissage de variance) et SVM ($C$, $\gamma$) ;
  \item \textbf{Random Search} (40 tirages) pour Random Forest et Gradient Boosting, dont
        l'espace (nombre d'arbres, profondeur, taux d'apprentissage, sous-échantillonnage\ldots)
        est trop grand pour une recherche exhaustive.
\end{itemize}
Le nombre de variables $k$ est optimisé conjointement. Le critère est l'AUC en validation
croisée stratifiée à 5 plis ; le pipeline complet (filtre, standardisation, sélection,
modèle) est réappris dans chaque pli. Enfin, un \textbf{stacking} des trois meilleurs modèles
optimisés est construit.

\capture{gridsearch_sortie}{Exécution des recherches d'hyperparamètres (sortie du notebook).}
  {Capture de la sortie de la cellule d'optimisation (section 4) : une ligne par modèle avec AUC, durée et meilleurs paramètres.}{5.5cm}

\subsection{Résultats}
\tabgen{tuning}{Gain apporté par l'optimisation (AUC en validation croisée).}{tuning}
\tabgen{params}{Meilleurs hyperparamètres trouvés.}{params}
\tabgen{cv_opt}{Validation croisée des modèles optimisés.}{cv_opt}
\figgen{avant_apres_optimisation}{F1 et AUC en validation croisée avant et après optimisation.}{1.0}

\subsection{Choix du modèle final}
On retient le modèle ayant la meilleure AUC en validation croisée ; en cas de quasi-égalité
(écart inférieur à 0,002) on préfère le modèle le plus simple puis celui dont l'écart
entraînement--validation est le plus faible, c'est-à-dire le meilleur compromis
biais/variance. Le modèle retenu est \textbf{\resfinalModel} (AUC en validation croisée
\resfinalCvAuc, F1 \resfinalCvFone, accuracy \resfinalCvAcc), avec une pénalité L1
($C=0{,}1$) et $k=20$ variables.

Après optimisation, presque tous les modèles atteignent une AUC de 1,000 en validation
croisée : le critère est \textbf{saturé} et ne permet plus de les départager. La règle de
départage choisit alors le modèle le plus simple, qui a aussi le plus petit écart
entraînement--validation et les meilleurs scores en validation à plis contigus
(tableau~\ref{tab:robustesse}). Ce choix est cohérent avec toutes les informations
disponibles avant le test.

% =====================================================================
\section{Évaluation finale sur l'échantillon de test}

Le modèle final est entraîné sur les 2000 premiers cycles puis évalué une seule fois sur les
cycles 2001 à 2205. Les autres modèles optimisés sont évalués à titre de comparaison
seulement : le choix a été fait avant, sur la validation croisée.

\begin{table}[H]\centering\small
\caption{Performances du modèle final (\resfinalModel) sur le test.}\label{tab:final}
\begin{tabular}{ccccc}
\toprule
Accuracy & Précision & Rappel & F1 & AUC\\
\midrule
\restestAccuracy & \restestPrecision & \restestRecall & \restestFone & \restestRocauc\\
\bottomrule
\end{tabular}
\quad
\begin{tabular}{lcc}
\toprule
& prédit non opt. & prédit optimale\\
\midrule
réel non optimale & \restestTN & \restestFP\\
réel optimale & \restestFN & \restestTP\\
\bottomrule
\end{tabular}
\end{table}

\tabgen{test}{Performances de tous les modèles optimisés sur le test (comparaison).}{test}
\figgen{test_final}{Matrice de confusion du modèle final et courbes ROC de tous les modèles sur le test.}{1.0}

Le modèle final ne fait \textbf{aucun faux négatif} (rappel 1) mais \resnTestErrors{} faux
positifs : des cycles à valve non optimale prédits optimaux. Son AUC sur le test vaut
pourtant \restestRocauc : les probabilités \emph{classent} parfaitement les cycles, mais le
seuil de 0,5 appris sur l'entraînement n'est plus adapté au test. À l'inverse, le SVM, la
forêt aléatoire, le gradient boosting et l'arbre de décision optimisés obtiennent 100\,\% sur
le test (tableau~\ref{tab:test}).

\subsection{Analyse des erreurs}
Le tableau~\ref{tab:erreurs_valve} détaille les erreurs selon le niveau réel de la valve et la
figure~\ref{fig:probas_test} montre la probabilité prédite pour chaque cycle de test.

\tabgen{erreurs_valve}{Erreurs du modèle final sur le test selon la condition réelle de la valve.}{erreurs_valve}
\figgen{probas_test}{Probabilité prédite d'une valve optimale pour chaque cycle de test ; les erreurs sont entourées.}{1.0}

\textbf{Toutes les erreurs concernent des valves à 90\,\%} (léger retard), soit le niveau le plus
proche de l'optimum ; les niveaux 80 et 73\,\% sont parfaitement reconnus. Sur la
figure~\ref{fig:probas_test}, ces cycles ont une probabilité entre 0,5 et 0,75 : ils sont mieux
classés que les vrais cycles optimaux (probabilité proche de 1) mais restent au-dessus du
seuil. L'explication est le décalage de distribution décrit en section 2 : dans le test,
69\,\% des cycles ont un accumulateur à 90\,bar (33\,\% à l'entraînement) et tous ont un
refroidisseur à 100\,\%. Ces conditions modifient légèrement le niveau moyen de la pression,
donc la valeur de \texttt{PS2\_w1\_mean} sur laquelle repose presque entièrement la régression
logistique, et la poussent juste au-dessus du seuil pour les valves à 90\,\%. Les modèles non
linéaires, qui découpent l'espace en plusieurs régions, sont moins sensibles à ce décalage.

\textbf{Leçon.} On ne change pas le modèle après avoir vu le test : ce serait ajuster le choix
sur le test, et le score annoncé ne serait plus une estimation honnête. Mais ce résultat
montre que, quand le critère de sélection est saturé, il faut un critère plus exigeant
(validation à plis contigus en F1, ou validation par groupes de conditions de
fonctionnement), et qu'en production un SVM ou une forêt aléatoire serait préférable, avec
une surveillance de la calibration du seuil.

% =====================================================================
\section{Interprétation : causes d'une valve non optimale}
L'importance par permutation (figure~\ref{fig:importance_permutation}) mesure la baisse
d'AUC sur le test quand on mélange une variable ; elle ne dépend pas du type de modèle.

\figgen{importance_permutation}{Importance par permutation des variables du modèle final (test).}{0.75}

Le modèle final repose presque entièrement sur une seule variable, \texttt{PS2\_w1\_mean}, la
pression moyenne entre 6 et 12\,s : la mélanger fait perdre 0,57 point d'AUC, alors que les
autres variables ont une importance quasi nulle (la pénalité L1 a annulé leurs coefficients
ou elles sont redondantes avec elle).

\paragraph{Interprétation physique.} Cette fenêtre contient la montée en pression de 9,4\,s.
Si la valve commute en retard, la pression monte plus tard, donc la pression moyenne sur la
fenêtre est plus faible. Pour le vérifier, on mesure directement l'instant où PS2 dépasse
50\,bar après 5\,s (figure~\ref{fig:instant_commutation}). Sur l'entraînement, cet instant vaut
en moyenne \resriseHundred\,s pour une valve à 100\,\%, \resriseNinety\,s à 90\,\%,
\resriseEighty\,s à 80\,\% et \resriseSeventythree\,s à 73\,\%, avec une dispersion de quelques
millisecondes seulement et \textbf{sans dépendance à l'état de l'accumulateur}.

\figgen{instant_commutation}{Instant où la pression PS2 dépasse 50 bar, selon la condition de la valve (entraînement et test).}{0.8}

\textbf{La cause d'une valve non optimale est donc un retard de commutation} : dès le niveau
90\,\%, la montée en pression arrive environ \resdelayNinetyMs\,ms plus tard, et ce retard
grandit avec la dégradation. C'est cette signature que tous les modèles exploitent,
directement (instant du maximum, raideur de la montée) ou indirectement (moyenne de pression
sur la fenêtre, énergie haute fréquence).

\paragraph{Expérience a posteriori.} Un classifieur logistique entraîné sur cette seule
variable obtient sur le test une accuracy de \resoneFeatTestAcc{} et un F1 de
\resoneFeatTestFone. Cette expérience a été menée après l'évaluation finale et ne remplace pas
le modèle choisi ; elle montre qu'une variable directement liée au mécanisme physique (un
instant, plutôt qu'un niveau de pression) est beaucoup plus robuste aux changements de
conditions de fonctionnement.

% =====================================================================
\section{Bonus réalisés}\label{sec:bonus}

\subsection{Tests unitaires}
Le dossier \texttt{tests/} contient 11 tests \texttt{pytest} qui vérifient, sur des signaux
simulés : la lecture des fichiers et la numérotation des cycles, la binarisation de la cible,
le découpage 2000\,/\,205 sans recouvrement, la valeur des caractéristiques sur un signal
connu, la somme des parts d'énergie spectrale (égale à 1), le filtre de corrélation, le bon
fonctionnement des pipelines et l'API de l'application web.
\begin{lstlisting}[language=bash]
python -m pytest -q
\end{lstlisting}
\capture{tests_pytest}{Exécution des tests unitaires.}
  {Capture du terminal après \texttt{python -m pytest -q} (« 11 passed »).}{4cm}

\subsection{Conteneurisation}
Le \texttt{Dockerfile} construit une image Python 3.11 contenant le code, les dépendances et le
modèle entraîné. Elle permet à un tiers d'exécuter le projet sans rien installer :
\begin{lstlisting}[language=bash]
docker build -t valve-ml .
docker run --rm -p 5000:5000 valve-ml                       # application web
docker run --rm valve-ml python -m pytest -q                 # tests
docker run --rm -v "$PWD/data:/app/data" -v "$PWD/rapport:/app/rapport" valve-ml \
  jupyter nbconvert --to notebook --execute notebooks/projet_valve.ipynb --output projet_valve.ipynb
\end{lstlisting}
\capture{docker}{Construction et exécution du conteneur Docker.}
  {Capture du terminal montrant \texttt{docker build} terminé puis \texttt{docker run} de l'application.}{5cm}

\subsection{Application web}
L'application Flask (\texttt{app/app.py}) charge le modèle final et renvoie, pour un numéro de
cycle donné, la prédiction, la probabilité associée, la condition réelle et si le cycle
appartient au jeu d'entraînement ou de test. Une API JSON est aussi disponible :
\texttt{GET /api/predict?cycle=2050}.
\capture{appli_web}{Application web : prédiction pour un cycle du jeu de test.}
  {Capture du navigateur sur \texttt{http://localhost:5000/?cycle=2050} (ou un autre cycle de test) montrant le résultat.}{7cm}

\subsection{Dépôt GitHub}
\capture{github}{Dépôt GitHub du projet.}
  {Capture de la page d'accueil du dépôt GitHub (arborescence et README). Lien du dépôt : [à compléter].}{6cm}

% =====================================================================
\section{Limites et perspectives}
\begin{itemize}
  \item \textbf{Décalage de distribution.} Le test correspond à une période où l'accumulateur
        est fortement dégradé et le refroidisseur en parfait état. Les performances sur le test
        doivent donc être lues comme une mesure de robustesse à des conditions peu vues à
        l'entraînement.
  \item \textbf{Critère de sélection saturé.} L'AUC en validation croisée vaut 1 pour presque
        tous les modèles ; la validation à plis contigus, ou par groupes de conditions de
        fonctionnement, devrait servir directement à l'optimisation et au choix du modèle.
  \item \textbf{Calibration du seuil.} L'AUC du modèle final sur le test est parfaite : un seuil
        recalibré périodiquement sur des cycles récents étiquetés corrigerait ses erreurs.
  \item \textbf{Caractéristiques physiques.} Ajouter explicitement l'instant de commutation (et
        celui du pic de débit) aux variables rendrait les modèles plus robustes.
  \item \textbf{Capteurs.} Seuls PS2 et FS1 sont autorisés ; d'autres capteurs du banc (PS1,
        EPS1, températures) pourraient améliorer la robustesse.
  \item \textbf{Approches sans extraction manuelle.} Un réseau convolutif 1D sur les signaux
        bruts, ou des caractéristiques automatiques (\emph{tsfresh}), sont des pistes
        naturelles.
  \item \textbf{Classification multi-classe.} Prédire le niveau exact (100, 90, 80, 73\,\%)
        permettrait de planifier la maintenance avant la panne.
\end{itemize}

% =====================================================================
\section{Conclusion}
Ce projet a mis en place une chaîne complète de maintenance prédictive : contrôle de qualité,
construction de la cible binaire, analyse exploratoire, extraction de \resnFeatures{}
caractéristiques physiques, sélection de variables sans fuite d'information, comparaison de
neuf modèles, évaluation par validation croisée stratifiée, analyse biais/variance et
optimisation des hyperparamètres par Grid Search et Random Search. Le modèle retenu,
\resfinalModel, obtient sur les \resnTest{} cycles de test une accuracy de \restestAccuracy, un
F1 de \restestFone{} et une AUC de \restestRocauc. Ses erreurs viennent d'un seuil décalé par
le changement de conditions de fonctionnement entre les périodes d'entraînement et de test ;
les modèles non linéaires optimisés atteignent 100\,\% sur le test.

Sur le plan métier, l'analyse répond à la question posée : une valve non optimale se
reconnaît à un \textbf{retard de la montée en pression} dans le cycle, mesurable dès un
léger retard de commutation (niveau 90\,\%) et indépendant de l'état des autres composants.
Surveiller cet instant de commutation constitue un indicateur simple et robuste de
maintenance prédictive.

% =====================================================================
\appendix
\section{Reproduire les résultats}
\begin{lstlisting}[language=bash]
pip install -r requirements.txt
# placer PS2.txt, FS1.txt et profile.txt dans data/
cd notebooks && jupyter nbconvert --to notebook --execute projet_valve.ipynb --output projet_valve.ipynb
cd ../rapport && pdflatex rapport.tex && pdflatex rapport.tex
python ../app/app.py
\end{lstlisting}

\section*{Référence}
N. Helwig, E. Pignanelli, A. Schütze, \emph{Condition monitoring of a complex hydraulic system
using multivariate statistics}, Proc. I2MTC-2015, IEEE International Instrumentation and
Measurement Technology Conference, 2015. Jeu de données : UCI Machine Learning Repository,
\emph{Condition monitoring of hydraulic systems}.

\end{document}
